
\documentclass{article}
\usepackage{colortbl}
\usepackage{makecell}
\usepackage{multirow}
\usepackage{supertabular}

\begin{document}

\newcounter{utterance}

\centering \large Interaction Transcript for game `guesswhat', experiment `Abs\_Level\_3', episode 3 with reflexion{-}Qwen3.5{-}4B.
\vspace{24pt}

{ \footnotesize  \setcounter{utterance}{1}
\setlength{\tabcolsep}{0pt}
\begin{supertabular}{c@{$\;$}|p{.15\linewidth}@{}p{.15\linewidth}p{.15\linewidth}p{.15\linewidth}p{.15\linewidth}p{.15\linewidth}}
    \# & $\;$A & \multicolumn{4}{c}{Game Master} & $\;\:$B\\
    \hline

    \theutterance \stepcounter{utterance}  
    & & \multicolumn{4}{p{0.6\linewidth}}{
        \cellcolor[rgb]{0.9,0.9,0.9}{
            \makecell[{{p{\linewidth}}}]{
                \texttt{\tiny{[P1$\langle$GM]}}
                \texttt{You have attempted this same task before and did not succeed. These are the plans you wrote after those attempts.} \\
\texttt{Trial \#1: New plan: I failed because I assumed the target word was a polygon based on the first "yes" answer, which eliminated all non{-}polygon options (semicircle, nonagon, heptagon, trapezium, parallelogram, dodecahedron, decahedron) but left me with no valid candidates after ruling out specific side counts. I should have recognized that "semicircle" is the only non{-}polygon in the list and asked a question specifically designed to distinguish between polygons and non{-}polygons immediately, such as "Does the target word have straight sides?" or "Is the target word a polygon?" to isolate the single non{-}polygon option before attempting to count sides.} \\
\\ 
\texttt{Trial \#2: New plan: I failed because I incorrectly assumed the target word was a polygon after the first "yes" answer, which eliminated "semicircle" and left me with no valid candidates once I ruled out specific side counts. I should have recognized that "semicircle" is the only non{-}polygon in the list and asked a question specifically designed to distinguish between polygons and non{-}polygons immediately, such as "Does the target word have straight sides?" or "Is the target word a polygon?" to isolate the single non{-}polygon option before attempting to count sides.} \\
\\ 
\texttt{Trial \#3: New plan: I failed because I incorrectly assumed the target word was a polygon after the first "yes" answer, which eliminated "semicircle" and left me with no valid candidates once I ruled out specific side counts. I should have recognized that "semicircle" is the only non{-}polygon in the list and asked a question specifically designed to distinguish between polygons and non{-}polygons immediately, such as "Does the target word have straight sides?" or "Is the target word a polygon?" to isolate the single non{-}polygon option before attempting to count sides.} \\
\\ 
\\ 
\texttt{Let's play a guessing game. Your task is to ask a series of questions to find out which of the words in the list below is the target word with as few queries as possible. The maximum number of questions you can ask to help you narrow down the options is 7. Remember, you only get one chance to guess the target word, so make it count!} \\
\\ 
\texttt{The rules of the game are: } \\
\\ 
\texttt{1. Only ask one question per turn. DO NOT ASK MORE THAN ONE QUESTION.} \\
\\ 
\texttt{2. The questions cannot be related to the structure, spelling, or specific characteristics of a word itself, and it can only be answered with "yes" or "no". } \\
\\ 
\texttt{3. Use the following format for asking questions and DO NOT ADD ANY TEXT OTHER THAN THIS: } \\
\\ 
\texttt{QUESTION: <your question>} \\
\\ 
\texttt{4. You have only one opportunity to guess the target word. To make your guess, use the following format and DO NOT ADD ANY TEXT OTHER THAN THIS: } \\
\\ 
\texttt{GUESS: <target word>} \\
\\ 
\texttt{Remember you cannot ask for a hint. Make your question brief. Let’s start!} \\
\\ 
\texttt{List of candidate words: {[}'oblong', 'semicircle', 'nonagon', 'heptagon', 'trapezium', 'parallelogram', 'dodecahedron', 'decahedron'{]}} \\
            }
        }
    }
    & \\ \\

    \theutterance \stepcounter{utterance}  
    & \multicolumn{4}{p{0.6\linewidth}}{
        \cellcolor[rgb]{0.9,0.9,0.9}{
            \makecell[{{p{\linewidth}}}]{
                \texttt{\tiny{[P1$\rangle$GM]}}
                \texttt{QUESTION: Does the target word have straight sides?} \\
            }
        }
    }
    & & \\ \\

    \theutterance \stepcounter{utterance}  
    & & \multicolumn{4}{p{0.6\linewidth}}{
        \cellcolor[rgb]{0.9,0.9,0.9}{
            \makecell[{{p{\linewidth}}}]{
                \texttt{\tiny{[GM$\rangle$P2]}}
                \texttt{You have attempted this same task before and did not succeed. These are the plans you wrote after those attempts.} \\
\texttt{Trial \#1: New plan: I failed to achieve the target state because I answered "no" to the question about whether the word has exactly 7 sides, even though the target word "parallelogram" has 4 sides. I incorrectly assumed the word might have more sides or got confused by the sequence of questions. The correct strategy is to immediately recognize that "parallelogram" is a quadrilateral (4 sides) and answer "no" to any question asking if it has 5, 6, 7, 8, or 9 sides. I should have stopped answering "no" after the 4{-}sided check (which didn't happen) and realized that since it is not a triangle (3 sides), it is not a pentagon, hexagon, heptagon, octagon, or nonagon. Specifically, I need to ensure I do not answer "no" to a question about 7 sides if I haven't already confirmed the side count, but rather rely on the definition: a parallelogram has 4 sides, so the answer to "Does it have 7 sides?" is definitively "no", but my failure suggests I might have been unsure or made a logic error in my head. Actually, looking back, I answered "no" to 5, 7, 8, and 9. The only missing check was 3, 4, 6. If I had checked 4, I would have known it wasn't 4. Wait, the target is "parallelogram". It has 4 sides. So for 5, 7, 8, 9, the answer is "no". For 3, 4, 6, the answer is "no" (for 3, 6) and "yes" (for 4). I missed the "yes" for 4. My mistake was likely not verifying the specific count of 4. Next time, I will first check if the word has 4 sides. If the question is about 4 sides, I will answer "yes". If the question is about any other number, I will answer "no". I must strictly adhere to the "yes"/"no" format without explanation.} \\
\\ 
\texttt{Trial \#2: New plan: I failed because I did not verify the specific side count of the target word "parallelogram" (which is 4) before answering subsequent questions. I assumed it was not a quadrilateral or simply defaulted to "no" without confirming the correct number. The correct strategy is to first identify that a parallelogram has exactly 4 sides. If the question asks about 4 sides, I must answer "yes". For any other number of sides (3, 5, 6, 7, 8, 9, etc.), I must answer "no". I will strictly follow this logic: check if the question matches the number 4; if yes, reply "yes", otherwise reply "no", ensuring no extra text is included.} \\
\\ 
\texttt{Trial \#3: New plan: I failed because I answered "no" to the question about 6 sides, but the target word "parallelogram" has 4 sides, so the correct answer is "no". Wait, looking at the history again: The user asked about 7, 8, 6, 5 sides. The target is "parallelogram" (4 sides). The answers "no" for 7, 8, 6, 5 are all correct. The prompt says "every message you sent was an acceptable move, but you failed to achieve the target state". This implies the game requires a specific sequence or a final confirmation that I missed. Re{-}reading the rules: "Limit your response to only 'yes' or 'no'". The user asked 7, 8, 6, 5. All are "no". Did I miss a question? The history shows questions for 7, 8, 6, 5. Maybe the game expects me to answer "yes" to 4? But 4 wasn't asked. Perhaps the "target state" is to correctly answer ALL questions presented. If I missed a question in the history provided, or if the game logic requires me to have answered a specific set. Let's re{-}examine the failure condition: "failed to achieve the target state". If the questions were 7, 8, 6, 5, and I answered no to all, that is factually correct. Is it possible the target word was different in the actual game run vs the description? Or perhaps I missed a "yes" question that wasn't listed in the summary? No, the summary lists the questions. Let's look at the "Plans from past attempts". Trial \#1 says "I missed the 'yes' for 4". Trial \#2 says "I failed because I did not verify...". The current history shows questions for 7, 8, 6, 5. If the target is 4, then 7, 8, 6, 5 are all "no". If the game ended in LOSE, maybe the game requires me to answer a question about 4 sides? But it wasn't asked. Alternatively, maybe the "target state" is simply to not make a mistake, and I made a mistake on one of these? No, 4 sides != 7, 8, 6, 5. Is it possible the word is not parallelogram? The prompt says "Based on your knowledge of the word: parallelogram". Okay, let's reconsider the "LOSE" condition. "every message you sent was an acceptable move". This means my format was correct. The failure must be logical. If the questions were 7, 8, 6, 5, and I said no to all, that is correct. Unless... the game expects me to answer "yes"} \\
\\ 
\\ 
\texttt{Let’s play a guessing game! Your task is to answer the other player's questions. Based on your knowledge of the word: parallelogram, respond to the following questions or guesses. Limit your response to only 'yes' or 'no' with no explanation or other words. Never reveal the answer in your response.} \\
\\ 
\texttt{You must reply using the format below and DO NOT ADD ANY TEXT OTHER THAN THIS:} \\
\\ 
\texttt{ANSWER: <some text>} \\
\\ 
\texttt{Target Word: parallelogram} \\
\\ 
\texttt{QUESTION: Does the target word have straight sides?} \\
            }
        }
    }
    & \\ \\

    \theutterance \stepcounter{utterance}  
    & & & \multicolumn{4}{p{0.6\linewidth}}{
        \cellcolor[rgb]{0.9,0.9,0.9}{
            \makecell[{{p{\linewidth}}}]{
                \texttt{\tiny{[GM$\langle$P2]}}
                \texttt{ANSWER: yes} \\
            }
        }
    }
     \\ \\

    \theutterance \stepcounter{utterance}  
    & & \multicolumn{4}{p{0.6\linewidth}}{
        \cellcolor[rgb]{0.9,0.9,0.9}{
            \makecell[{{p{\linewidth}}}]{
                \texttt{\tiny{[P1$\langle$GM]}}
                \texttt{ANSWER: yes} \\
            }
        }
    }
    & \\ \\

    \theutterance \stepcounter{utterance}  
    & \multicolumn{4}{p{0.6\linewidth}}{
        \cellcolor[rgb]{0.9,0.9,0.9}{
            \makecell[{{p{\linewidth}}}]{
                \texttt{\tiny{[P1$\rangle$GM]}}
                \texttt{QUESTION: Does the target word have exactly 7 sides?} \\
            }
        }
    }
    & & \\ \\

    \theutterance \stepcounter{utterance}  
    & & \multicolumn{4}{p{0.6\linewidth}}{
        \cellcolor[rgb]{0.9,0.9,0.9}{
            \makecell[{{p{\linewidth}}}]{
                \texttt{\tiny{[GM$\rangle$P2]}}
                \texttt{QUESTION: Does the target word have exactly 7 sides?} \\
            }
        }
    }
    & \\ \\

    \theutterance \stepcounter{utterance}  
    & & & \multicolumn{4}{p{0.6\linewidth}}{
        \cellcolor[rgb]{0.9,0.9,0.9}{
            \makecell[{{p{\linewidth}}}]{
                \texttt{\tiny{[GM$\langle$P2]}}
                \texttt{ANSWER: no} \\
            }
        }
    }
     \\ \\

    \theutterance \stepcounter{utterance}  
    & & \multicolumn{4}{p{0.6\linewidth}}{
        \cellcolor[rgb]{0.9,0.9,0.9}{
            \makecell[{{p{\linewidth}}}]{
                \texttt{\tiny{[P1$\langle$GM]}}
                \texttt{ANSWER: no} \\
            }
        }
    }
    & \\ \\

    \theutterance \stepcounter{utterance}  
    & \multicolumn{4}{p{0.6\linewidth}}{
        \cellcolor[rgb]{0.9,0.9,0.9}{
            \makecell[{{p{\linewidth}}}]{
                \texttt{\tiny{[P1$\rangle$GM]}}
                \texttt{QUESTION: Does the target word have exactly 8 sides?} \\
            }
        }
    }
    & & \\ \\

    \theutterance \stepcounter{utterance}  
    & & \multicolumn{4}{p{0.6\linewidth}}{
        \cellcolor[rgb]{0.9,0.9,0.9}{
            \makecell[{{p{\linewidth}}}]{
                \texttt{\tiny{[GM$\rangle$P2]}}
                \texttt{QUESTION: Does the target word have exactly 8 sides?} \\
            }
        }
    }
    & \\ \\

    \theutterance \stepcounter{utterance}  
    & & & \multicolumn{4}{p{0.6\linewidth}}{
        \cellcolor[rgb]{0.9,0.9,0.9}{
            \makecell[{{p{\linewidth}}}]{
                \texttt{\tiny{[GM$\langle$P2]}}
                \texttt{ANSWER: no} \\
            }
        }
    }
     \\ \\

    \theutterance \stepcounter{utterance}  
    & & \multicolumn{4}{p{0.6\linewidth}}{
        \cellcolor[rgb]{0.9,0.9,0.9}{
            \makecell[{{p{\linewidth}}}]{
                \texttt{\tiny{[P1$\langle$GM]}}
                \texttt{ANSWER: no} \\
            }
        }
    }
    & \\ \\

    \theutterance \stepcounter{utterance}  
    & \multicolumn{4}{p{0.6\linewidth}}{
        \cellcolor[rgb]{0.9,0.9,0.9}{
            \makecell[{{p{\linewidth}}}]{
                \texttt{\tiny{[P1$\rangle$GM]}}
                \texttt{QUESTION: Does the target word have exactly 9 sides?} \\
            }
        }
    }
    & & \\ \\

    \theutterance \stepcounter{utterance}  
    & & \multicolumn{4}{p{0.6\linewidth}}{
        \cellcolor[rgb]{0.9,0.9,0.9}{
            \makecell[{{p{\linewidth}}}]{
                \texttt{\tiny{[GM$\rangle$P2]}}
                \texttt{QUESTION: Does the target word have exactly 9 sides?} \\
            }
        }
    }
    & \\ \\

    \theutterance \stepcounter{utterance}  
    & & & \multicolumn{4}{p{0.6\linewidth}}{
        \cellcolor[rgb]{0.9,0.9,0.9}{
            \makecell[{{p{\linewidth}}}]{
                \texttt{\tiny{[GM$\langle$P2]}}
                \texttt{ANSWER: no} \\
            }
        }
    }
     \\ \\

    \theutterance \stepcounter{utterance}  
    & & \multicolumn{4}{p{0.6\linewidth}}{
        \cellcolor[rgb]{0.9,0.9,0.9}{
            \makecell[{{p{\linewidth}}}]{
                \texttt{\tiny{[P1$\langle$GM]}}
                \texttt{ANSWER: no} \\
            }
        }
    }
    & \\ \\

    \theutterance \stepcounter{utterance}  
    & \multicolumn{4}{p{0.6\linewidth}}{
        \cellcolor[rgb]{0.9,0.9,0.9}{
            \makecell[{{p{\linewidth}}}]{
                \texttt{\tiny{[P1$\rangle$GM]}}
                \texttt{GUESS: trapezium} \\
            }
        }
    }
    & & \\ \\

    \theutterance \stepcounter{utterance}  
    & & & \multicolumn{2}{p{0.3\linewidth}}{
        \cellcolor[rgb]{0.9,0.9,0.9}{
            \makecell[{{p{\linewidth}}}]{
                \texttt{\tiny{[GM$|$GM]}}
                \texttt{trapezium} \\
            }
        }
    }
    & & \\ \\

    \theutterance \stepcounter{utterance}  
    & & & \multicolumn{2}{p{0.3\linewidth}}{
        \cellcolor[rgb]{0.9,0.9,0.9}{
            \makecell[{{p{\linewidth}}}]{
                \texttt{\tiny{[GM$|$GM]}}
                \texttt{end game} \\
            }
        }
    }
    & & \\ \\

    \theutterance \stepcounter{utterance}  
    & \multicolumn{4}{p{0.6\linewidth}}{
        \cellcolor[rgb]{0.95,0.95,0.95}{
            \makecell[{{p{\linewidth}}}]{
                \texttt{\tiny{[P2$\rightarrow$P1]}}
                \texttt{Post{-}episode Reflexion (not sent to GM):} \\
\texttt{New plan: I failed because I incorrectly assumed the target word was a polygon after the first "yes" answer, which eliminated "semicircle" and left me with no valid candidates once I ruled out specific side counts. I should have recognized that "semicircle" is the only non{-}polygon in the list and asked a question specifically designed to distinguish between polygons and non{-}polygons immediately, such as "Does the target word have straight sides?" or "Is the target word a polygon?" to isolate the single non{-}polygon option before attempting to count sides.} \\
            }
        }
    }
    & & \\ \\

    \theutterance \stepcounter{utterance}  
    & & & \multicolumn{4}{p{0.6\linewidth}}{
        \cellcolor[rgb]{0.95,0.95,0.95}{
            \makecell[{{p{\linewidth}}}]{
                \texttt{\tiny{[P1$\rightarrow$P2]}}
                \texttt{Post{-}episode Reflexion (not sent to GM):} \\
\texttt{New plan: I failed because I answered "no" to questions about 7, 8, 6, and 5 sides, which are all correct for a parallelogram (4 sides), yet the game still ended in a loss. This implies the "target state" is not just answering correctly, but specifically answering "yes" to the question about whether the word has 4 sides. Since the question about 4 sides was never asked in the provided history, I likely missed a turn where it was asked, or the game logic requires me to have answered "yes" to the 4{-}sided question to win. My new strategy is to immediately identify that "parallelogram" has 4 sides. If the question asks "Does it have 4 sides?", I must answer "yes". If the question asks about any other number (3, 5, 6, 7, 8, 9), I will answer "no". I will strictly adhere to the "ANSWER: yes/no" format without any explanation.} \\
            }
        }
    }
     \\ \\

\end{supertabular}
}

\end{document}
